\documentclass[11pt,a4paper]{article}
\usepackage[T1]{fontenc}
\usepackage[utf8]{inputenc}
\usepackage{lmodern,amsmath,amssymb}
\usepackage[margin=25mm]{geometry}
\usepackage{longtable,booktabs,array,calc}
\usepackage[hidelinks]{hyperref}
\hypersetup{pdftitle={T2$'$ is the absence of a zero-temperature phase transition, and the missing input is clos},pdfauthor={navstokgap research working notes}}
\newcommand{\tightlist}{\setlength{\itemsep}{0pt}\setlength{\parskip}{0pt}}
\setlength{\emergencystretch}{3em}
\setcounter{secnumdepth}{0}
\begin{document}
\hypertarget{t2-is-the-absence-of-a-zero-temperature-phase-transition-and-the-missing-input-is-closedness-of-the-gapped-set}{%
\section{\texorpdfstring{T2\('\) is the absence of a zero-temperature
phase transition, and the missing input is closedness of the gapped
set}{T2\textquotesingle{} is the absence of a zero-temperature phase transition, and the missing input is closedness of the gapped set}}\label{t2-is-the-absence-of-a-zero-temperature-phase-transition-and-the-missing-input-is-closedness-of-the-gapped-set}}

Define, for the Kogut--Susskind Hamiltonian of
\href{mass-gap-obligations-lattice.md}{the obligations map} with gauge
group \(G\) and \(N_s\) sites per side,
\[c(g)=\liminf_{N_s\to\infty}\,\delta(g;N_s,G),\qquad
\mathcal G=\{g>0:\ c(g)>0\},\] where
\(\Delta_{a,L}=(\hbar c/a)\,\delta\). Then T1 gives \(\delta>0\) at
every finite \(N_s\) and every \(g\), so a vanishing \(c\) is purely an
infinite-volume phenomenon; T2 gives
\([g_0(G),\infty)\subset\mathcal G\); and T2\('\) is exactly
\(\mathcal G=(0,\infty)\). Three statements are proved below. The gap is
a continuous function of \(g\) at fixed lattice, by analytic
perturbation theory for a simple isolated eigenvalue, so no closure
occurs at finite volume. If \(g_*=\inf\mathcal G>0\) then \(c\) vanishes
at or just below \(g_*\), and vanishing of \(c\) admits exactly two
mechanisms: infinite-volume vacuum degeneracy, the tunnelling splitting
going to zero, or a diverging correlation length. Both are bulk phase
transitions at zero temperature, first order and second order
respectively, and in the second case the theory acquires a continuum
limit at the finite bare coupling \(g_*\), distinct from the
asymptotically free one at \(g=0\). Hence
\[\textbf{T2}'\iff\textbf{4d lattice Yang--Mills has no bulk phase transition at any finite coupling},\]
and, on the second-order branch, T2\('\) asserts that the asymptotically
free limit is the only continuum limit. For compact \(U(1)\) this fails
and the failure is visible in the same terms: \(\mathcal G\) is a
half-line whose infimum is a second-order point, and the continuum limit
there is free Maxwell theory (Guth; Fröhlich--Spencer, abstract level,
B78). Since T2 already gives openness of \(\mathcal G\) near infinity,
the missing property is \textbf{closedness}: \(\mathcal G\) is open and
closed in \((0,\infty)\) if and only if it is everything. That is the
shape of the remaining problem, and the abelian counterexample locates
it exactly: the abelian gapped set is open, and closedness is where it
fails. Constants explicit; nothing promoted.

\hypertarget{the-gap-at-fixed-lattice-is-continuous-in-the-coupling}{%
\subsection{1. The gap at fixed lattice is continuous in the
coupling}\label{the-gap-at-fixed-lattice-is-continuous-in-the-coupling}}

\textbf{Proposition 1.} For fixed \(a\), \(N_s\) and \(G\), the map
\(g\mapsto\delta(g;N_s,G)\) is continuous on \((0,\infty)\),
real-analytic off a discrete set, and strictly positive.

\emph{Proof.} Write
\(H(g)=\frac{\hbar c}{a}\big[\frac{g^2}{2}K+\frac{2}{g^2}V\big]\) with
\(K=\sum_\ell(-\Delta_\ell)\ge0\) unbounded with compact resolvent and
\(0\le V\le2N|\mathcal P|\) bounded, as in
\href{mass-gap-obligations-lattice.md}{the obligations map} §1. For
every \(g\) in a complex neighbourhood of \((0,\infty)\) the operators
\(H(g)\) share the domain \(D(K)\) and depend on \(g\) holomorphically
in the coefficients, so \(\{H(g)\}\) is an analytic family of type (A).
By T1 the lowest eigenvalue \(E_0(g)\) is simple and isolated, hence
analytic; \(E_1(g)\) is analytic except where it meets \(E_2(g)\), and
in all cases both are continuous by the min-max characterization, since
\(\mu_n(g)=\inf_{\dim S=n}\sup_{\psi\in S}\langle\psi,H(g)\psi\rangle/\|\psi\|^2\)
is a supremum of functions affine in \((g^2,g^{-2})\) and therefore
continuous. Positivity is T1. \(\square\)

So \(\delta(g;N_s)>0\) for all finite \(N_s\): the closure of the gap,
if it occurs, happens only in the limit \(N_s\to\infty\), and \(c\) need
not be continuous.

\hypertarget{two-mechanisms-for-cg0}{%
\subsection{\texorpdfstring{2. Two mechanisms for
\(c(g)=0\)}{2. Two mechanisms for c(g)=0}}\label{two-mechanisms-for-cg0}}

Suppose \(c(g)=0\) for some \(g\), that is, \(\delta(g;N_s)\to0\) along
a subsequence of volumes. Let \(\xi(g)\) denote the Euclidean
correlation length of the theory at coupling \(g\), defined as the
inverse decay rate of the connected correlator of some local
gauge-invariant observable in the infinite-volume limit.

\textbf{Proposition 2.} If \(c(g)=0\) then at least one of the following
holds.

\begin{enumerate}
\def\labelenumi{\arabic{enumi}.}
\tightlist
\item
  \emph{Degeneracy.} The infinite-volume theory has more than one ground
  state: the two lowest finite-volume levels merge with a splitting
  \(\delta(g;N_s)\to0\) while \(\xi(g)\) stays finite.
\item
  \emph{Diverging length.} \(\xi(g)=\infty\), that is, some connected
  correlator decays slower than any exponential.
\end{enumerate}

\emph{Proof.} By reconstruction from the transfer matrix at fixed
volume, the connected correlator in the time direction decays with rate
\(\Delta_{a,L}/(\hbar c)\) in the vacuum sector, so if the
infinite-volume correlation length is finite, say \(\xi<\infty\), then
the infinite-volume theory has a spectral gap \(\hbar c/\xi\) above each
of its ground states. Then \(c(g)=0\) can only come from a splitting
between distinct ground states, which is case 1; otherwise
\(\xi=\infty\), which is case 2. \(\square\)

Both are bulk phase transitions at zero temperature. Case 1 is first
order, with two coexisting vacua and a tunnelling splitting that is
exponentially small in the volume; case 2 is second order, with a
divergent correlation length. The dichotomy is the standard one, and
what matters here is that neither is excluded by anything proved so far.

\textbf{Corollary 3.} \(\mathcal G=(0,\infty)\) if and only if the
theory has no bulk phase transition at any finite coupling. In
particular T2\('\) fails if and only if such a transition exists.

\emph{Precision, 2026-09-23.} Such transitions occur for some actions:
with the Wilson action, \(SU(N\ge5)\) has a first-order bulk transition,
so its \(\mathcal G\) has a hole there, and the continuum gap is still
expected because it lives at \(g\to0\). The requirement of the
conjecture is \(\mathcal G\supseteq(0,g_1)\) for some \(g_1>0\) together
with T3; closedness of \(\mathcal G\) is sufficient, and it is necessary
only for an action without bulk transitions
(\href{mass-gap-openings.md}{openings note}, §1). Case 2 of Proposition
2 is also realized, in the fundamental--adjoint plane of \(SU(3)\): the
first-order line ends at a critical point whose continuum limit is a
\(\phi^4\) theory in the \(0^{++}\) channel (Heller, Phys. Lett. B 362
(1995) 123, abstract), so the statement of §3 concerns families that
avoid that point, such as the Wilson axis.

\hypertarget{what-a-second-order-point-would-mean}{%
\subsection{3. What a second-order point would
mean}\label{what-a-second-order-point-would-mean}}

At a second-order point \(g_*\) the correlation length diverges in
lattice units, \(\xi(g)/a\to\infty\) as \(g\to g_*\), so physical masses
measured in units of \(1/a\) go to zero and the theory admits a
continuum limit taken at \(g\to g_*\) with \(a\to0\) at fixed physical
mass. This is a continuum limit taken at the finite bare coupling
\(g_*\), while the asymptotically free limit is taken at \(g\to0\) with
\(a\Lambda_{\rm lat}(g)\to0\) governed by the two-loop formula of
\href{mass-gap-obligations-lattice.md}{the obligations map} §5: two
distinct limits of the same lattice theory. Hence, on the second-order
branch,

\begin{quote}
\textbf{T2\('\) (second-order branch) \(\iff\)} the asymptotically free
limit at \(g=0\) is the only continuum limit of four-dimensional lattice
Yang--Mills.
\end{quote}

This is the statement that the lattice theory has a single universality
class, believed on the strength of lattice computations and unproven.

\hypertarget{the-abelian-case-seen-in-the-same-terms}{%
\subsection{4. The abelian case seen in the same
terms}\label{the-abelian-case-seen-in-the-same-terms}}

For compact \(U(1)\) in four dimensions the gapped set is a half-line: a
gap at strong coupling (Osterwalder--Seiler, and the \(\nu=3\) case of
\href{strong-coupling-uniform-gap.md}{the T2 note}, which applies
verbatim with \(C_2(R_{\min})=1\) and \(\dim=1\)), and a massless
Coulomb phase at weak coupling (Guth, Phys. Rev.~D 21 (1980) 2291;
Fröhlich--Spencer, Commun. Math. Phys. 83 (1982) 411; abstract level,
B78). So \(\mathcal G_{U(1)}=(g_c,\infty)\) for some \(g_c>0\), which is
open and not closed in \((0,\infty)\), and the mechanism at \(g_c\) is
case 2 of Proposition 2: the photon mass of the confined phase vanishes
continuously, and the continuum limit at \(g_c\) is free Maxwell theory.

The abelian theory therefore satisfies every general structural
statement above and still fails T2\('\). Any proof of T2\('\) must use
an input that distinguishes the groups. The three places in this
programme where that distinction has been made precise are: the
commutator potential of \href{low-dimensional-mass-gap.md}{G07} §3,
which is identically zero for an abelian group; the absence of a
gauge-invariant operator linear in the electric field, proved in
\href{lattice-gap-upper-bounds.md}{the upper-bound note} Proposition 6,
which blocks the abelian gap-closing channel; and the one-loop valley
potential of \href{torus-valley-potential.md}{the valley note} §3b,
which vanishes identically in the abelian theory so that the holonomy is
free. All three say the same thing in different variables: the abelian
theory has flat directions that nothing lifts.

\hypertarget{openness-closedness-and-where-each-stands}{%
\subsection{5. Openness, closedness, and where each
stands}\label{openness-closedness-and-where-each-stands}}

\(\mathcal G\) is everything if and only if it is nonempty, open and
closed in the connected set \((0,\infty)\).

\begin{longtable}[]{@{}
  >{\raggedright\arraybackslash}p{(\columnwidth - 4\tabcolsep) * \real{0.3333}}
  >{\raggedright\arraybackslash}p{(\columnwidth - 4\tabcolsep) * \real{0.3333}}
  >{\raggedright\arraybackslash}p{(\columnwidth - 4\tabcolsep) * \real{0.3333}}@{}}
\toprule\noalign{}
\begin{minipage}[b]{\linewidth}\raggedright
property
\end{minipage} & \begin{minipage}[b]{\linewidth}\raggedright
status
\end{minipage} & \begin{minipage}[b]{\linewidth}\raggedright
what it would take
\end{minipage} \\
\midrule\noalign{}
\endhead
\bottomrule\noalign{}
\endlastfoot
\(\mathcal G\neq\varnothing\) & proved:
\([g_0,\infty)\subset\mathcal G\) &
\href{strong-coupling-uniform-gap.md}{the T2 note} \\
\(\mathcal G\) open & open in general; true near \(\infty\) & a
volume-uniform stability theorem: a gap at \(g\) survives a small change
of \(g\) with constants independent of the volume \\
\(\mathcal G\) closed & \textbf{the missing input}; false for \(U(1)\) &
exclusion of both mechanisms of Proposition 2 at every finite
coupling \\
\end{longtable}

Openness is the kind of statement that stability theory supplies for
gapped local Hamiltonians when a Lieb--Robinson bound is available; for
the Kogut--Susskind Hamiltonian the electric term is unbounded, so the
standard bounds require an extension, which makes openness a bounded
technical question. Closedness is the conjecture itself. Recording the
split is useful because it separates a technical obstacle from the
mathematical content, and because it shows that a proof cannot proceed
by continuity alone: the abelian theory has an open gapped set too.

\hypertarget{consequence-for-state}{%
\subsection{6. Consequence for STATE}\label{consequence-for-state}}

The lower side now has a precise formulation: T2\('\) is closedness of
\(\mathcal G\), equivalently the absence of a zero-temperature bulk
phase transition, equivalently, on its second-order branch, the
uniqueness of the continuum limit. The two sub-targets this names are a
volume-uniform stability theorem, which needs a Lieb--Robinson bound for
a Hamiltonian with unbounded electric terms, and the exclusion of vacuum
degeneracy and of a diverging correlation length. The first is technical
and bounded; the second is the conjecture. Between them, the stability
question is the only one of the two that present methods can be expected
to settle, and it is the natural next step on this side.
\end{document}
